\documentclass[a4paper,11pt]{article}
\usepackage[T1]{fontenc}
\usepackage[utf8]{inputenc}
\usepackage[italian]{babel}
\usepackage{amsmath,amssymb}
\usepackage{booktabs}
\usepackage{longtable}
\usepackage{geometry}
\usepackage{hyperref}
\usepackage{xcolor}
\usepackage{enumitem}
\usepackage{graphicx}

\geometry{margin=2.5cm}

\title{Dai descrittori al filtro bicomb\\[0.5em]
\large Strategia di mappatura per \emph{interfantasia}}
\author{Francesco Ferrante}
\date{Maggio 2026}

\begin{document}
\maketitle

\tableofcontents
\newpage

% ============================================================
\section{Inquadramento del problema}
% ============================================================

Il sistema dispone di sedici descrittori spettrali (centroid, spread,
rolloff, slope, flatness, crest, skewness, kurtosis, tonality, tpr, flux,
irregularity, zcr, entropy, obsir\_std, zscore) calcolati frame per frame
sul segnale del musicista. Il filtro bicomb dispone di cinque ingressi:

\begin{itemize}[nosep]
  \item tre parametri continui a rampa, ciascuno nell'intervallo
        $[-1, +1]$, che possono assumere qualsiasi valore intermedio;
  \item due contatori a evento, che si incrementano di $+1$ al passaggio
        di una condizione e governano rispettivamente la lunghezza del
        ramo FIR e del ramo IIR.
\end{itemize}

La domanda di ricerca è: come collegare i sedici descrittori a questi
cinque ingressi in modo che il sistema sia musicalmente leggibile, sia
per chi suona sia per chi ascolta?

% ============================================================
\section{Strategie classiche e loro limiti}
% ============================================================

La letteratura su feature extraction in live electronics (Bullock 2008,
Jehan 2005, Park 2004) offre quattro famiglie di strategie:

\begin{enumerate}
  \item Vettore di feature in input a un classificatore (k-NN, ANN)
        per riconoscere classi sonore.
  \item Combinazione lineare pesata di poche feature scelte, scalata e
        filtrata, per ottenere un controllo continuo.
  \item Soglie sui valori o sulle derivate temporali delle feature per
        generare eventi discreti.
  \item Uso di feature singole come segnali di controllo simbolici (f0,
        onset, envelope).
\end{enumerate}

Tutte queste strategie chiedono di scegliere a priori quali descrittori
combinare. La scelta è di solito ad hoc rispetto al brano e al pezzo
musicale. Per il nostro caso, in cui lo stesso impianto deve funzionare
su un corpus eterogeneo, abbiamo cercato un criterio più sistematico.

% ============================================================
\section{Da un'analisi di correlazione a una scelta}
% ============================================================

Prima di decidere il mapping abbiamo studiato la matrice di correlazione
dei sedici descrittori sul corpus, in due regimi.

\paragraph{Correlazione fra-gesti.}
Aggregando tutti i frame di tutti i gesti, emerge un cluster molto coeso
(centroid, rolloff, spread, flatness, irregularity, entropy, zcr, crest)
con $\lvert r\rvert > 0{,}7$ stabilmente fra strumenti diversi.
Tonality e tpr appaiono anti-correlati a questo cluster. Solo skewness e
kurtosis formano una coppia indipendente più debole. Slope, flux,
obsir\_std e zscore mostrano comportamenti variabili fra cataloghi.

\paragraph{Correlazione intra-gesto.}
Calcolando invece la correlazione \emph{dentro} il singolo gesto
(cioè quanto i descrittori si muovono insieme istante per istante),
la struttura collassa: solo la coppia flatness/tonality resta correlata
(media $-0{,}99$, deviazione $0{,}01$ su 257 gesti). Tutte le altre coppie
mostrano grande variabilità: lo stesso paio di descrittori può essere
fortemente correlato in un gesto e scorrelato in un altro.

\paragraph{Conseguenza.}
La struttura di correlazione fra descrittori \emph{non è una costante
strutturale del sistema}, ma varia con il segnale. Questo è un dato
musicalmente prezioso: significa che la matrice di correlazione, calcolata
in tempo reale, è essa stessa una sorgente di informazione sul carattere
del suono, indipendente dai valori dei singoli descrittori.

% ============================================================
\section{La strategia adottata: la matrice come sorgente}
% ============================================================

La strategia proposta è la seguente:

\begin{enumerate}
  \item A ogni istante $t$ si calcola la matrice di correlazione di
        Pearson $M(t)$ fra i sedici descrittori, su una finestra mobile
        di $N$ frame.
  \item Dalla matrice si estraggono cinque scalari (tre proprietà
        strutturali, due indicatori di evoluzione) che diventano
        direttamente i cinque ingressi del bicomb.
\end{enumerate}

Il vantaggio concettuale è doppio. Primo, partecipa l'intera tabella:
nessuna casella viene scelta a priori, e nessun descrittore viene
preferito. Secondo, i controlli misurano la \emph{grammatica} con cui il
timbro evolve, non i valori istantanei dello spettro. Il filtro bicomb
risponde al modo in cui il timbro si organizza, non al timbro in sé.

% ============================================================
\section{Calcolo della matrice rolling}
% ============================================================

A ogni frame $t \ge N - 1$, sia $\mathbf{x}_i(t)$ la serie temporale del
descrittore $i$ ristretta alla finestra $[t - N + 1,\, t]$. La cella
$(i, j)$ della matrice $M(t)$ è il coefficiente di correlazione di
Pearson fra $\mathbf{x}_i(t)$ e $\mathbf{x}_j(t)$:
\[
M_{ij}(t) = \frac{\sum_{k} (x_i^{(k)} - \bar x_i)(x_j^{(k)} - \bar x_j)}
                 {\sqrt{\sum_{k} (x_i^{(k)} - \bar x_i)^2}\,
                  \sqrt{\sum_{k} (x_j^{(k)} - \bar x_j)^2}}.
\]

La matrice è simmetrica, ha diagonale costante $+1$ e contiene
$K(K-1)/2$ valori distinti nel triangolo superiore (per $K = 15$
descrittori utilizzabili, sono 105 numeri). Tutti i valori stanno in
$[-1, +1]$, il che si rivela cruciale: i controlli che ne derivano sono
già nella scala richiesta dal filtro.

Sui descrittori con varianza nulla nella finestra (caso tipico delle
sinusoidi pure) la correlazione è indefinita; la cella corrispondente
viene marcata e ignorata nelle aggregazioni successive.

% ============================================================
\section{I cinque controlli}
% ============================================================

% --------------------------------------------------------
\subsection{$C_1$ -- Coerenza media segnata}
% --------------------------------------------------------

\formula
\[
C_1(t) = \frac{2}{K(K-1)} \sum_{i < j} M_{ij}(t).
\]

\paragraph{Significato.}
Media di tutte le correlazioni a coppie con il loro segno. Misura il
livello medio di sincronia fra i descrittori nella finestra corrente.

\paragraph{Interpretazione musicale.}
\begin{itemize}[nosep]
  \item $C_1 \to +1$: tutti i descrittori si muovono nello stesso verso.
        Tipico dei gesti in cui un inviluppo molto marcato trascina con sé
        ogni aspetto del timbro (colpo percussivo, crescendo armonico).
  \item $C_1 \to 0$: i descrittori si muovono indipendentemente.
        Tipico dei suoni stazionari (rumore, sinusoide tenuta) o dei
        gesti articolati dove diverse dimensioni dello spettro
        cambiano in modo non coordinato.
  \item $C_1 \to -1$: anti-correlazione dominante. Caso raro, sintomo
        di trasformazioni molto strutturate (es. crescendo che scurisce).
\end{itemize}

\paragraph{Range osservato.}
Sul corpus, $C_1$ si muove in $[-0{,}05,\, +0{,}30]$. I valori più alti
sono prodotti dai gesti del clarinetto contrabbasso ($C_{1,\text{media}}
\approx +0{,}12$) e dalle sinusoidi in crescendo. Sui suoni stazionari
(rumore bianco, sinusoide a 440 Hz) e sui timpani tenuti, $C_1$ resta
vicino allo zero.

% --------------------------------------------------------
\subsection{$C_2$ -- Concentrazione assiale}
% --------------------------------------------------------

\formula
Siano $\lambda_1 \ge \lambda_2 \ge \dots \ge \lambda_K$ gli autovalori
in modulo della matrice $M(t)$, ordinati in modo decrescente. Allora:
\[
C_2(t) = 2\,\frac{\lvert \lambda_1 \rvert}{\sum_{i=1}^{K} \lvert \lambda_i \rvert} - 1.
\]

\paragraph{Significato.}
Frazione dell'informazione totale della matrice catturata dal primo asse
principale, riscalata da $[0, 1]$ a $[-1, +1]$. È una misura di
\emph{dimensionalità effettiva} del sistema descrittori.

\paragraph{Interpretazione musicale.}
\begin{itemize}[nosep]
  \item $C_2 \to +1$: la matrice è quasi rank-1. Tutto il movimento dei
        sedici descrittori è proiettabile su un solo asse. Il suono è
        \emph{monolitico}: tutte le sue dimensioni variano in modo
        coordinato.
  \item $C_2 \to -1$: gli autovalori sono distribuiti uniformemente.
        Il sistema occupa molte dimensioni indipendenti, e il suono ha
        una struttura interna ricca, non riducibile a un singolo
        ``percorso''.
\end{itemize}

\paragraph{Range osservato.}
$C_2$ si muove in $[-0{,}45,\, +0{,}40]$ a seconda dei gesti. I suoni
sintetici stazionari (sinusoidi, rumori filtrati) hanno valori medi
positivi (intorno a $+0{,}17$, ${+0{,}20}$). I gesti con distorsione
non-lineare evolutiva (\texttt{tanh\_drive}) producono escursioni ampie
da $-0{,}40$ a $+0{,}26$, segno che la dimensionalità del sistema cambia
durante il gesto.

% --------------------------------------------------------
\subsection{$C_3$ -- Asimmetria di segno}
% --------------------------------------------------------

\formula
Siano $n_+$ il numero di celle del triangolo superiore con $M_{ij} > 0$
e $n_-$ il numero di celle con $M_{ij} < 0$. Allora:
\[
C_3(t) = \frac{n_+ - n_-}{n_+ + n_-}.
\]

\paragraph{Significato.}
Bilanciamento netto fra coppie concordi e coppie discordi. A differenza
di $C_1$, non considera la magnitudine: due matrici possono avere stessa
media $C_1$ ma topologia di segni diversa.

\paragraph{Interpretazione musicale.}
\begin{itemize}[nosep]
  \item $C_3 \to +1$: dominano le concordanze. Suoni con organizzazione
        coerente, anche se debole.
  \item $C_3 \to -1$: dominano le opposizioni. Caso raro nel corpus.
  \item $C_3 \to 0$: equilibrio fra concordanze e opposizioni.
\end{itemize}

\paragraph{Range osservato.}
$C_3$ si muove in $[-0{,}10,\, +0{,}35]$. Anche qui il clarinetto si
distingue (media $+0{,}17$). Il valore di $C_3$ tende a seguire $C_1$
ma non in modo deterministico: nei gesti articolati le due grandezze
si discostano.

% --------------------------------------------------------
\subsection{$E_1$ -- Velocità di evoluzione strutturale}
% --------------------------------------------------------

\formula
Sia $D$ il ritardo in frame. Si definisce la distanza di Frobenius
normalizzata fra la matrice attuale e quella di $D$ frame fa:
\[
E_1(t) = \frac{1}{\sqrt{K(K-1)}}\,
        \Bigl\| M(t) - M(t - D) \Bigr\|_F
      = \frac{1}{\sqrt{K(K-1)}}\,
        \sqrt{\sum_{i,j} \bigl( M_{ij}(t) - M_{ij}(t-D) \bigr)^2}.
\]

La normalizzazione per $\sqrt{K(K-1)}$ rende $E_1$ indipendente dalla
taglia della matrice: una soglia fissata su $E_1$ ha lo stesso
significato qualunque sia il numero di descrittori.

\paragraph{Significato.}
Tasso a cui la struttura delle relazioni fra descrittori si trasforma.

\paragraph{Trigger.}
$E_1$ alimenta uno dei due contatori a evento del bicomb. La condizione
di trigger è:
\[
E_1(t-1) \le \theta_{E_1} < E_1(t),
\]
con un tempo di refrattarietà di $\sim 200$ ms dopo ogni trigger per
evitare conteggi multipli sullo stesso evento.

\paragraph{Range e soglia.}
Su finestra $N = 32$, $D = 16$, $E_1$ si muove in $[0, 0{,}8]$ circa.
La soglia di default $\theta_{E_1} = 0{,}15$ produce mediamente 2--3
trigger per gesto. Sul timpano, dove gli inviluppi sono lenti ma
articolati, salgono a $\sim 2{,}6$ trigger; sui sintetici stazionari
restano vicini a $3{,}5$ per via di micro-instabilità delle correlazioni.

% --------------------------------------------------------
\subsection{$E_2$ -- Rotazione dell'asse dominante}
% --------------------------------------------------------

\formula
Sia $\mathbf{v}_1(t)$ il primo autovettore (associato a $\lambda_1$)
della matrice $M(t)$. Si definisce:
\[
E_2(t) = \frac{\mathbf{v}_1(t) \cdot \mathbf{v}_1(t - D)}
              {\lVert\mathbf{v}_1(t)\rVert\, \lVert\mathbf{v}_1(t-D)\rVert}.
\]

Il segno dell'autovettore è arbitrario (numericamente); per renderlo
deterministico si fissa positivo il primo elemento non nullo. La
quantità di interesse è il valore assoluto del coseno.

\paragraph{Significato.}
Coseno dell'angolo fra l'asse principale corrente e quello di $D$ frame
fa. Vale $+1$ quando l'orientamento dell'asse non è cambiato, scende
verso $0$ quando l'asse è ruotato di $90^\circ$.

\paragraph{Trigger.}
$E_2$ alimenta l'altro contatore. La condizione è opposta a $E_1$
(trigger \emph{in discesa}, cioè quando il coseno cala sotto soglia):
\[
\lvert E_2(t-1) \rvert \ge \theta_{E_2} > \lvert E_2(t) \rvert.
\]

\paragraph{Range e soglia.}
Sui gesti stazionari $E_2$ resta sopra $0{,}95$; sui gesti articolati può
scendere a $0{,}4$ e oltre. Con $\theta_{E_2} = 0{,}9$, sul corpus si
ottengono mediamente tre-quattro trigger per gesto; i timpani
producono il numero massimo ($\sim 6{,}2$ trigger di media), coerente
con suoni che ruotano lentamente il proprio asse spettrale durante il
decay.

% ============================================================
\section{Parametri di taratura}
% ============================================================

I parametri da fissare sono:

\begin{itemize}[nosep]
  \item \textbf{$N$ -- finestra rolling per la correlazione.} Troppo
        corta ($N < 16$) produce correlazioni molto rumorose; troppo
        lunga ($N > 128$) rende il sistema lento e perde le dinamiche
        rapide. Default attuale: $N = 32$ frame, pari a circa
        $0{,}75$ s con \texttt{hop} di $\sim 23$ ms.
  \item \textbf{$D$ -- ritardo per $E_1$ ed $E_2$.} Tipicamente $D = N/2$.
        Default: $D = 16$ frame, circa $0{,}37$ s.
  \item \textbf{$\theta_{E_1}$ -- soglia di trigger sull'evoluzione.}
        Default $0{,}15$; abbassandola si ottengono trigger più frequenti.
  \item \textbf{$\theta_{E_2}$ -- soglia di trigger sulla rotazione.}
        Default $0{,}9$.
  \item \textbf{Refrattarietà.} Tempo morto fra trigger successivi.
        Default $200$ ms.
\end{itemize}

% ============================================================
\section{Caratterizzazione del corpus}
% ============================================================

L'analisi batch sui 256 gesti del corpus (sette cataloghi: clarinetto
contrabbasso, timpano tenuto, sintetici, test\_segnali a $-30$ dB, tre
serie \texttt{recs}-002/003/004) produce le seguenti caratteristiche
medie per catalogo:

\begin{center}
\begin{tabular}{lccccc}
\toprule
catalogo & $C_1$ media & $C_1$ range & $C_2$ media & $E_1$ max & trigger $E_2$\\
\midrule
clarinetto       & $+0{,}12$ & $0{,}36$ & $+0{,}05$ & $0{,}73$ & $4{,}0$\\
timpano          & $+0{,}01$ & $0{,}17$ & $-0{,}05$ & $0{,}63$ & $6{,}2$\\
sintetici        & $-0{,}02$ & $0{,}11$ & $+0{,}17$ & $0{,}36$ & $3{,}1$\\
test\_segnali    & $-0{,}03$ & $0{,}11$ & $+0{,}20$ & $0{,}35$ & $3{,}1$\\
recs-002         & $-0{,}00$ & $0{,}10$ & $-0{,}04$ & $0{,}48$ & $3{,}2$\\
recs-003         & $+0{,}02$ & $0{,}12$ & $-0{,}04$ & $0{,}48$ & $3{,}2$\\
recs-004         & $+0{,}03$ & $0{,}11$ & $-0{,}02$ & $0{,}47$ & $3{,}1$\\
\bottomrule
\end{tabular}
\end{center}

\paragraph{Osservazioni principali.}

\begin{itemize}
  \item Il clarinetto si distingue su $C_1$ ed escursione di $C_1$,
        confermando che i suoi gesti hanno inviluppi che trascinano
        l'intera struttura dei descrittori.
  \item Il timpano è il catalogo più attivo su $E_2$: i suoi suoni
        tenuti evolvono lentamente nell'orientamento spettrale, anche
        senza grandi variazioni di intensità.
  \item I sintetici hanno $C_2$ medio positivo, segno di matrici più
        concentrate sull'asse principale (suoni a struttura semplice).
  \item Le tre serie \texttt{recs} si comportano in modo quasi identico
        fra loro su tutte le metriche, coerentemente con la loro natura
        di repliche dello stesso protocollo di test.
\end{itemize}

% ============================================================
\section{Questioni aperte}
% ============================================================

Tre punti restano da decidere sperimentalmente prima dell'implementazione
in Pure Data.

\paragraph{Indipendenza di $C_3$ rispetto a $C_1$.}
Sul corpus le due grandezze tendono a co-variare. Andrà verificato se il
guadagno informativo di $C_3$ giustifica la sua presenza come terzo
ingresso indipendente, oppure se conviene sostituirlo con un'altra
grandezza scalare ricavata dalla matrice (per esempio l'entropia della
distribuzione degli autovalori, o la kurtosis della distribuzione delle
correlazioni nel triangolo).

\paragraph{Filtraggio a valle.}
I controlli $C_1$, $C_2$, $C_3$ vanno smussati prima di alimentare il
filtro: una media mobile leggera (8 frame) o un EMA con $\alpha \approx
0{,}1$ è sufficiente. Per i trigger $E_1$ e $E_2$ è opportuno applicare
una \emph{soft dead zone} attorno alla soglia con isteresi.

\paragraph{Pre-condizionamento del segnale in tempo reale.}
In Pd il segnale del microfono deve essere passato attraverso un gate e
un compressore prima dell'estrazione dei descrittori, come prescritto
in Bullock (2008) per Sparkling Box. I valori del gate sono già fissati
per l'analisi offline ($-65$ dBFS assoluto, $-30$ dB relativo al picco
globale).

% ============================================================
\section{Architettura per famiglie e mappatura definitiva}
% ============================================================

L'esplorazione sulla matrice globale $15 \times 15$ ha mostrato che da essa
si estraggono al massimo due controlli continui indipendenti ($C_1$ e
$C_2$): ogni altra grandezza derivata ($C_3$, $\sigma_M$ della
distribuzione delle correlazioni, kurtosis, polarizzazione, frazione di
celle deboli) si rivela un doppione strutturale di uno dei due. La
matrice globale, da sola, non porta abbastanza dimensioni indipendenti
per coprire i cinque ingressi del bicomb in modo non ridondante.

\subsection{Suddivisione in famiglie}

I sedici descrittori si raggruppano in quattro famiglie, ciascuna
dedicata a un aspetto del timbro:

\begin{center}
\begin{tabular}{lll}
\toprule
famiglia & $K_F$ & descrittori \\
\midrule
collocazione & 4 & centroid, rolloff, spread, slope \\
grana   & 5 & flatness, tonality, tpr, irregularity, entropy \\
forma        & 3 & skewness, kurtosis, crest \\
dinamica     & 3 & flux, zcr, obsir\_std (zscore opzionale) \\
\bottomrule
\end{tabular}
\end{center}

Per ogni famiglia si calcola una matrice di correlazione rolling
indipendente $M_F(t)$ di taglia $K_F \times K_F$, con finestra adattata:
$N_F = \max(16,\, 4 K_F)$, $D_F = N_F/2$. Da ognuna si estraggono i
quattro scalari $C_1^F$, $C_2^F$, $E_1^F$, $E_2^F$ con le stesse formule
della matrice globale.

\paragraph{Verifica empirica.}
Sul corpus dei 256 gesti, i $C_1^F$ delle quattro famiglie risultano
quasi ortogonali fra loro (correlazione cross-gesto massima $0{,}39$
fra collocazione e forma, il resto sotto $0{,}22$), contro lo $0{,}92$
fra $C_1$ e $C_3$ della matrice globale. Su un brano reale
(\emph{A Pierre senza elettronica}, due strumenti distinti registrati
sui due canali), i profili medi dei $C_1^F$ separano nettamente
flauto contrabbasso e clarinetto contrabbasso, e dentro ciascun
canale i controlli coprono il range pieno $[-1, +1]$. La struttura per
famiglie è quindi sia indipendente sia espressiva.

\subsection{Mappatura definitiva al bicomb}

Il bicomb riceve cinque ingressi: due controlli continui (frequenza di
taglio $f_c$ e fattore $Q$), un selettore discreto a quattro forme, e
due contatori a evento (lunghezza del ramo FIR e del ramo IIR). La
mappatura definitiva assegna a ciascun ingresso una sorgente specifica
fra i sedici scalari delle famiglie:

\begin{center}
\begin{tabular}{lll}
\toprule
ingresso bicomb & sorgente & famiglia guida \\
\midrule
continuo 1: $f_c$         & $C_1^{\text{collocazione}}$ & collocazione \\
continuo 2: $Q$           & $C_1^{\text{grana}}$   & grana \\
selettore a 4 forme       & $\arg\max_F C_1^F$ con isteresi & tutte e quattro \\
trigger FIR               & $E_1^{\text{dinamica}}$     & dinamica \\
trigger IIR               & $E_2^{\text{forma}}$        & forma \\
\bottomrule
\end{tabular}
\end{center}

Ogni famiglia ha un ruolo nominato: collocazione e grana pilotano
i due controlli continui (i due assi del filtro), forma e dinamica
producono gli eventi che allungano le catene FIR e IIR, e tutte e
quattro partecipano collettivamente al selettore di forma.

\paragraph{Lettura musicale.}
I due continui misurano \emph{com'è fatto} il suono in quel momento
(quanto luminoso, quanto rumoroso). I due trigger misurano
\emph{cosa sta succedendo} al suono (la dinamica produce eventi quando
la trama temporale si riassetta, la forma produce eventi quando il
profilo della distribuzione spettrale ruota). Il selettore dice quale
famiglia sta tirando il timbro adesso, e quindi che geometria deve
avere il filtro.

\subsection{Semantica dei controlli continui: integratore senza ritorno}

I due controlli continui $C_1^{\text{collocazione}}$ e
$C_1^{\text{grana}}$ \emph{non} sono valori di posizione assoluti
da assegnare a $f_c$ e $Q$. Sono \emph{spinte}, cioè velocità di
spostamento. Il bicomb ha uno stato interno (la $f_c$ e la $Q$ correnti)
che si aggiorna integrando la spinta:
\[
f_c(t) = f_c(t-1) + \alpha_{f}\,C_1^{\text{collocazione}}(t),
\qquad
Q(t) = Q(t-1) + \alpha_{Q}\,C_1^{\text{grana}}(t),
\]
con $\alpha_f, \alpha_Q$ coefficienti di taratura che fissano la
``velocità di reazione'' del filtro, e $f_c, Q$ saturati ai loro
limiti fisici.

\paragraph{Conseguenze.}
\begin{itemize}[nosep]
  \item \textbf{Lo zero non è uno stato stabile.} $C_1^F = 0$ significa
        spinta nulla: il filtro \emph{resta dove è stato lasciato}, non
        torna a un centro. I passaggi neutri del gesto non disturbano
        la geometria conquistata.
  \item \textbf{Memoria del gesto.} La posizione attuale di $f_c$ e $Q$
        è il risultato di tutta la storia delle spinte ricevute. Il
        musicista può ``portare'' il filtro in una zona spingendo nella
        direzione voluta, e poi lasciarlo lì interrompendo la
        coordinazione.
  \item \textbf{Dead-zone soft.} Per evitare drift dovuti al rumore di
        stima di $C_1^F$, vale la pena di applicare una dead-zone
        morbida di $\pm 0{,}03$--$0{,}05$ attorno allo zero. Non serve
        per stabilità (il sistema è già stabile a zero), ma per pulizia.
  \item \textbf{Convenzione di segno.} Va fissata una sola volta: i due
        valori positivi di $C_1^F$ aprono il filtro in una direzione,
        quelli negativi nell'altra. La scelta è compositiva e non
        cambia il calcolo offline.
\end{itemize}

\paragraph{Immagine intuitiva.}
Il filtro è un galleggiante con remi: galleggia dove sta, il musicista
può remare (coordinando collocazione o grana) per spostarlo,
smettere di remare lo lascia fermo, e ci sono due sponde che non si
possono attraversare (i limiti di $f_c$ e $Q$).

\subsection{Semantica dei trigger: eventi a soglia con refrattarietà}

Diversamente dai due controlli continui, $E_1^{\text{dinamica}}$ ed
$E_2^{\text{forma}}$ \emph{non} pilotano un'integrazione di stato. Sono
sorgenti di \emph{eventi discreti}: a ogni passaggio di soglia
incrementano di $+1$ il rispettivo contatore del bicomb (lunghezza del
ramo FIR per il primo, del ramo IIR per il secondo).

\paragraph{Condizioni di trigger.}
\begin{itemize}[nosep]
  \item $E_1^{\text{dinamica}}$ scatta \emph{in salita}, quando
        $E_1^{\text{dinamica}}(t-1) \le \theta_{E_1} < E_1^{\text{dinamica}}(t)$:
        cioè quando la struttura della famiglia dinamica si è
        riassettata abbastanza in fretta.
  \item $E_2^{\text{forma}}$ scatta \emph{in discesa},
        quando $|E_2^{\text{forma}}(t-1)| \ge \theta_{E_2} > |E_2^{\text{forma}}(t)|$:
        cioè quando il coseno fra autovettori successivi della famiglia
        forma cala sotto soglia, segno che l'asse principale della
        distribuzione spettrale è ruotato.
\end{itemize}

\paragraph{Refrattarietà.}
Dopo ogni trigger si applica un tempo morto di $\sim 200$ ms durante il
quale ulteriori passaggi di soglia vengono ignorati. Serve a evitare
conteggi multipli sullo stesso evento musicale (oscillazioni della
grandezza attorno alla soglia).

\paragraph{Effetto sul filtro.}
Ogni trigger \emph{allunga} la catena corrispondente del bicomb di un
tap. La lunghezza dei rami è quindi un altro tipo di memoria del
gesto: cresce nel tempo, contando gli eventi musicalmente rilevanti
prodotti dal musicista. Va deciso se prevedere un meccanismo di
azzeramento (a comando, oppure per decadimento lento), per evitare che
dopo una performance lunga i rami siano cresciuti troppo per essere
gestibili.

\subsection{Conseguenze sulle scelte precedenti}

Questa architettura per famiglie rende obsolete diverse scelte fatte
con la sola matrice globale:

\begin{itemize}[nosep]
  \item I controlli $C_1$, $C_2$ e $C_3$ \emph{globali} descritti nelle
        sezioni precedenti restano utili come strumenti di analisi e
        confronto, ma \emph{non sono} la sorgente dei controlli del
        bicomb. La matrice globale resta come riferimento, non come
        pipeline operativa.
  \item I trigger $E_1$ ed $E_2$ \emph{globali} sono sostituiti dalle
        loro versioni di famiglia ($E_1^{\text{dinamica}}$ e
        $E_2^{\text{forma}}$). Il calcolo è identico ma su una matrice
        più piccola, e gli eventi diventano localizzati alla famiglia
        che li produce.
  \item Le \emph{questioni aperte} sul terzo controllo continuo
        (sostituzione di $C_3$) sono chiuse: il terzo ingresso del
        bicomb non è un controllo continuo, è un selettore a cinque
        stati pilotato da $\arg\max_F C_1^F$.
\end{itemize}

% ============================================================
\section{Appendice: lettura dettagliata delle formule}
% ============================================================

Questa appendice riprende ogni formula della dispensa e ne descrive il
significato di ciascun simbolo, una componente alla volta. È pensata per
poter essere consultata senza dover ricostruire il contesto.

% --------------------------------------------------------
\subsection{Coefficiente di correlazione di Pearson}
% --------------------------------------------------------

\[
r = \frac{\sum_{i=1}^{n} (x_i - \bar x)(y_i - \bar y)}
         {\sqrt{\sum_{i=1}^{n} (x_i - \bar x)^2}\,
          \sqrt{\sum_{i=1}^{n} (y_i - \bar y)^2}}
\]

\paragraph{$n$.} Numero di campioni confrontati. Nel nostro caso è il
numero di frame della finestra rolling ($n = N$).

\paragraph{$x_i$, $y_i$.} Valori delle due serie temporali al frame $i$.
Tipicamente $x_i$ è il valore di un descrittore (per esempio il
centroide) al frame $i$, e $y_i$ il valore di un altro descrittore (per
esempio lo spread) allo stesso frame.

\paragraph{$\bar x$, $\bar y$.} Medie aritmetiche delle due serie sulla
finestra di $n$ frame:
\[
\bar x = \frac{1}{n}\sum_{i=1}^{n} x_i, \qquad
\bar y = \frac{1}{n}\sum_{i=1}^{n} y_i.
\]
Servono come punto di riferimento: si guarda quanto ogni frame si
discosta dalla media, non i valori assoluti.

\paragraph{$(x_i - \bar x)$, $(y_i - \bar y)$.} Scarti dalla media,
detti \emph{deviazioni centrate}. Sono positivi quando il frame è sopra
la media, negativi quando è sotto. Centrare significa traslare l'origine
sulla media: dopo questa operazione contano solo le oscillazioni.

\paragraph{$(x_i - \bar x)(y_i - \bar y)$.} Prodotto degli scarti al
frame $i$. Il segno dice se a quel frame i due descrittori si muovono
dalla stessa parte (prodotto positivo) o in direzioni opposte (prodotto
negativo). La grandezza pesa l'ampiezza congiunta della deviazione.

\paragraph{$\sum_{i=1}^{n} (x_i - \bar x)(y_i - \bar y)$.} Somma dei
prodotti su tutti i frame. È proporzionale alla \emph{covarianza}.
Misura, complessivamente, quanto le due serie si muovono in modo
concorde sulla finestra. Dipende dalle unità di misura: per renderla
adimensionale serve il denominatore.

\paragraph{$(x_i - \bar x)^2$.} Scarto al quadrato. Toglie il segno e
tiene solo la grandezza della deviazione.

\paragraph{$\sum_{i=1}^{n} (x_i - \bar x)^2$.} Somma dei quadrati degli
scarti. Misura quanto la serie $x$ varia complessivamente sulla finestra.
Se $x$ è quasi costante è piccola, se oscilla molto è grande. È
proporzionale alla varianza di $x$.

\paragraph{$\sqrt{\sum_{i=1}^{n} (x_i - \bar x)^2}$.} Radice della somma
dei quadrati. È proporzionale alla deviazione standard di $x$.
Rappresenta l'\emph{ampiezza} della serie centrata.

\paragraph{$\sqrt{\sum(x_i-\bar x)^2}\cdot\sqrt{\sum(y_i-\bar y)^2}$.}
Prodotto delle due ampiezze. Per la disuguaglianza di Cauchy--Schwarz è
il massimo valore (in modulo) che il numeratore può assumere. Dividere
per questa quantità normalizza $r$ nell'intervallo $[-1, +1]$ e toglie
ogni dipendenza dalle unità di misura.

\paragraph{Lettura geometrica.} Detti $\vec u$ e $\vec v$ i vettori
centrati $(x_i - \bar x)_{i=1\dots n}$ e $(y_i - \bar y)_{i=1\dots n}$,
il numeratore è il prodotto scalare $\vec u \cdot \vec v$, ciascun
fattore al denominatore è la norma $\lVert\vec u\rVert$ e
$\lVert\vec v\rVert$, quindi
\[
r = \frac{\vec u \cdot \vec v}{\lVert\vec u\rVert\,\lVert\vec v\rVert}
  = \cos\theta,
\]
dove $\theta$ è l'angolo fra le due serie centrate in $\mathbb{R}^n$.

% --------------------------------------------------------
\subsection{Matrice di correlazione rolling $M_{ij}(t)$}
% --------------------------------------------------------

\[
M_{ij}(t) = \frac{\sum_{k} (x_i^{(k)} - \bar x_i)(x_j^{(k)} - \bar x_j)}
                 {\sqrt{\sum_{k} (x_i^{(k)} - \bar x_i)^2}\,
                  \sqrt{\sum_{k} (x_j^{(k)} - \bar x_j)^2}}
\]

\paragraph{$t$.} Frame corrente. La matrice $M(t)$ è ricalcolata a ogni
frame, scorrendo la finestra di $N$ campioni.

\paragraph{$i$, $j$.} Indici dei due descrittori che si stanno
correlando. Con $K$ descrittori utilizzabili, $i, j \in \{1, \dots, K\}$.
La cella $(i, j)$ contiene la correlazione fra il descrittore $i$ e il
descrittore $j$.

\paragraph{$k$.} Indice del frame interno alla finestra rolling. Varia
fra $t - N + 1$ e $t$. Sostituisce l'indice $i$ del Pearson generico per
non confondersi con l'indice $i$ del descrittore.

\paragraph{$x_i^{(k)}$.} Valore del descrittore $i$ al frame $k$ della
finestra.

\paragraph{$\bar x_i$.} Media del descrittore $i$ sulla finestra di $N$
frame:
\[
\bar x_i = \frac{1}{N}\sum_{k=t-N+1}^{t} x_i^{(k)}.
\]

\paragraph{Struttura della matrice.} $M(t)$ è simmetrica
($M_{ij} = M_{ji}$), ha diagonale costante $M_{ii} = +1$ e contiene
$K(K-1)/2$ valori distinti nel triangolo superiore. Tutti i valori
stanno in $[-1, +1]$.

% --------------------------------------------------------
\subsection{$C_1$ -- Coerenza media segnata}
% --------------------------------------------------------

\[
C_1(t) = \frac{2}{K(K-1)} \sum_{i < j} M_{ij}(t)
\]

\paragraph{$K$.} Numero di descrittori utilizzabili (tipicamente 15 dopo
aver eliminato il descrittore degenerato per ogni frame).

\paragraph{$K(K-1)/2$.} Numero di coppie distinte $(i, j)$ con $i < j$.
È il numero di celle del triangolo superiore stretto della matrice.
Per $K = 15$ vale $105$.

\paragraph{$2/[K(K-1)]$.} Reciproco di $K(K-1)/2$, cioè $1/$(numero di
coppie). Trasforma la somma in una media.

\paragraph{$\sum_{i < j} M_{ij}(t)$.} Somma di tutte le correlazioni a
coppie nel triangolo superiore della matrice, con il loro segno
(concordanze positive aumentano, discordanze diminuiscono).

\paragraph{Esito.} $C_1 \in [-1, +1]$. Vale $+1$ solo se tutte le coppie
sono perfettamente correlate, $-1$ solo se tutte sono perfettamente
anti-correlate (caso non realizzabile per $K > 2$).

% --------------------------------------------------------
\subsection{$C_2$ -- Concentrazione assiale}
% --------------------------------------------------------

\[
C_2(t) = 2\,\frac{\lvert \lambda_1 \rvert}
                 {\sum_{i=1}^{K} \lvert \lambda_i \rvert} - 1
\]

\paragraph{$\lambda_1, \dots, \lambda_K$.} Autovalori della matrice
$M(t)$, ordinati in modulo decrescente:
$|\lambda_1| \ge |\lambda_2| \ge \dots \ge |\lambda_K|$. Essendo $M(t)$
simmetrica, sono tutti reali. Ogni autovalore corrisponde a una
direzione (autovettore) lungo cui la nuvola dei descrittori si estende:
$|\lambda_i|$ è proporzionale all'estensione su quella direzione.

\paragraph{$\sum_{i=1}^{K} |\lambda_i|$.} Somma dei moduli di tutti gli
autovalori. Equivale alla traccia di $M$ se la matrice è
semidefinita positiva (per Pearson, $\mathrm{tr}(M) = K$); usiamo i
moduli per gestire eventuali piccoli autovalori negativi prodotti da
rumore numerico.

\paragraph{$|\lambda_1| / \sum_i |\lambda_i|$.} Frazione di
``informazione'' (varianza) catturata dal primo asse principale. Sta in
$[1/K, 1]$: vale $1/K$ se gli autovalori sono tutti uguali (matrice
isotropa), vale $1$ se solo $\lambda_1$ è non nullo (matrice rank-1).

\paragraph{$2(\cdot) - 1$.} Riscalamento affine da $[0, 1]$ a $[-1, +1]$
per adattarsi alla scala del bicomb. In pratica si comprime ulteriormente
da $[1/K, 1]$ ma per $K = 15$ il minimo teorico ($-0{,}87$) non si
osserva mai sul corpus reale.

% --------------------------------------------------------
\subsection{$C_3$ -- Asimmetria di segno}
% --------------------------------------------------------

\[
C_3(t) = \frac{n_+ - n_-}{n_+ + n_-}
\]

\paragraph{$n_+$.} Numero di celle del triangolo superiore stretto di
$M(t)$ con $M_{ij} > 0$.

\paragraph{$n_-$.} Numero di celle con $M_{ij} < 0$. Le eventuali celle
con $M_{ij} = 0$ esatto (caso degenere) non contribuiscono.

\paragraph{$n_+ + n_-$.} Numero totale di celle considerate. Salvo
degenerazioni vale $K(K-1)/2$.

\paragraph{$n_+ - n_-$.} Eccesso di concordanze sulle opposizioni. Può
essere positivo o negativo.

\paragraph{Esito.} $C_3 \in [-1, +1]$. Vale $+1$ se tutte le coppie sono
positive, $-1$ se tutte negative. A differenza di $C_1$, ignora la
magnitudine delle correlazioni e guarda solo la loro topologia di segni.

% --------------------------------------------------------
\subsection{$E_1$ -- Distanza di Frobenius normalizzata}
% --------------------------------------------------------

\[
E_1(t) = \frac{1}{\sqrt{K(K-1)}}\,
        \sqrt{\sum_{i,j} \bigl( M_{ij}(t) - M_{ij}(t-D) \bigr)^2}
\]

\paragraph{$D$.} Ritardo in frame fra le due matrici confrontate.
Tipicamente $D = N/2$ (default: 16 frame).

\paragraph{$M_{ij}(t) - M_{ij}(t-D)$.} Variazione della singola cella
fra l'istante attuale e quello di $D$ frame fa. È un numero in
$[-2, +2]$.

\paragraph{$\bigl( M_{ij}(t) - M_{ij}(t-D) \bigr)^2$.} Variazione al
quadrato. Toglie il segno e amplifica le variazioni grandi rispetto a
quelle piccole.

\paragraph{$\sum_{i,j} (\dots)^2$.} Somma su tutte le celle (l'intera
matrice, non solo il triangolo). I termini sulla diagonale sono nulli
($M_{ii} = 1$ sempre), quindi contribuiscono solo i $K(K-1)$ termini
fuori diagonale, contati due volte per simmetria.

\paragraph{$\sqrt{\sum_{i,j} (\dots)^2}$.} Norma di Frobenius di
$M(t) - M(t-D)$, indicata $\lVert M(t) - M(t-D) \rVert_F$. È la distanza
euclidea fra le matrici lette come vettori in $\mathbb{R}^{K\times K}$.

\paragraph{$1/\sqrt{K(K-1)}$.} Fattore di normalizzazione. Senza di lui
la norma cresce con la taglia della matrice e una soglia $\theta_{E_1}$
dovrebbe essere riferita al numero di descrittori. Con questo fattore
$E_1$ è (in pratica) confinato in $[0, \sim 1]$ indipendentemente da $K$.

% --------------------------------------------------------
\subsection{$E_2$ -- Rotazione dell'asse dominante}
% --------------------------------------------------------

\[
E_2(t) = \frac{\mathbf{v}_1(t) \cdot \mathbf{v}_1(t - D)}
              {\lVert\mathbf{v}_1(t)\rVert\, \lVert\mathbf{v}_1(t-D)\rVert}
\]

\paragraph{$\mathbf{v}_1(t)$.} Primo autovettore di $M(t)$, cioè
l'autovettore associato all'autovalore di modulo maggiore $\lambda_1$.
È un vettore di $K$ componenti, una per descrittore. Indica la
\emph{direzione} dello spazio dei descrittori lungo cui la nuvola di
punti si estende di più nella finestra corrente.

\paragraph{$\mathbf{v}_1(t-D)$.} Stessa cosa, calcolata $D$ frame prima.

\paragraph{$\mathbf{v}_1(t) \cdot \mathbf{v}_1(t-D)$.} Prodotto scalare
fra i due autovettori. Misura quanto le due direzioni puntano nello
stesso verso.

\paragraph{$\lVert\mathbf{v}_1(t)\rVert$.} Norma euclidea
dell'autovettore. Le routine numeriche restituiscono autovettori già
normalizzati a 1, quindi questo fattore vale $1$ in pratica; resta nella
formula per chiarezza.

\paragraph{Rapporto.} È il coseno dell'angolo fra le due direzioni
principali. Vale $+1$ se l'asse non è ruotato, $0$ se è ruotato di
$90^\circ$, $-1$ se è invertito. Il segno è arbitrario: numericamente si
fissa positivo il primo elemento non nullo, e si lavora sul valore
assoluto.

\paragraph{Lettura musicale.} Quando l'asse dominante ruota, vuol dire
che il \emph{profilo} di quale combinazione di descrittori sta variando
di più è cambiato: per esempio, prima del cambio il timbro variava
soprattutto sull'asse luminosità (centroid, rolloff), dopo varia
soprattutto sull'asse grana (flatness, entropia). $E_2$ basso
segnala questo riassetto qualitativo del gesto.

% ============================================================
\section{Riferimenti}
% ============================================================

\begin{itemize}[nosep]
  \item Bullock, J. (2008). \emph{Implementing audio feature extraction
        in live electronic music}. PhD thesis, Birmingham City University.
  \item Peeters, G. (2003). \emph{A large set of audio features for sound
        description}. CUIDADO Project Report, IRCAM.
  \item Lerch, A. (2023). \emph{An Introduction to Audio Content Analysis}.
        Wiley.
  \item Park, T.\,H. (2004). \emph{Towards Automatic Musical Instrument
        Timbre Recognition}. PhD thesis, Princeton University.
\end{itemize}

\end{document}
